\section{Super-heavies}

\subsection{Ascalon}

Mounted on a Falchion chassis, the Ascalon keeps that tank's resilience but changes its role entirely. The Ascalon uses an Inferno Cannon, a devastating weapon that can flood its opponents with a torrent of flame. It is used to close on the most entrenched fortifications and scour them with purifying fire.

The Ascalon has the \textbf{Wounds (3)} and \textbf{DR (5+)} abilities. Its Ascalon Inferno Cannon uses a \textbf{Template (Flamer)} and has the \textbf{Reduces Cover ($-3$)} ability.

\begin{unitstatstable}
\SetCell[r=3]{c,m} 15 cm &
\SetCell[r=3]{c,m} 2+ &
\SetCell[r=3]{c,m} +6 &
\SetCell[r=3]{c,m} 4 &
Ascalon Inferno Cannon &
0 cm &
Template &
4+ &
$-1$
\\
& & & &
Twin Heavy Bolters &
45 cm &
1 &
4+ &
$-1$
\\
& & & &
Laser Destroyer &
75 cm &
2 &
4+ &
$-3$
\end{unitstatstable}

\unitimage{ascalon.png}

\unitdivider

\subsection{Cerberus}

The Cerberus heavy tank is a super-heavy tank hunter that was used in limited numbers by the Space Marine Legions during the Great Crusade. The Cerberus is a rare and experimental variant of the super-heavy assault tank. Both use the same hull, armour, and chassis design. The Cerberus was designed by the Mechanicum of Mars in order to field-test archeotech weapons recovered during the Great Crusade, known as Neutron Projectors. The Cerberus is an example of one of the many war machines that were at first used only by a few Legions and, given their results, became common among all the Legions.

The Cerberus has the \textbf{Wounds (2)} and \textbf{DR (6+)} abilities. Its Twin Neutron Projector inflicts \textbf{Multiple Hits (1D3)}.

\begin{unitstatstable}
\SetCell[r=2]{c,m} 15 cm &
\SetCell[r=2]{c,m} 2+ &
\SetCell[r=2]{c,m} +5 &
\SetCell[r=2]{c,m} 4 &
Twin Neutron Projector &
75 cm &
1 &
3+ &
$-4$
\\
& & & &
Lascannon &
60 cm &
2 &
5+ &
$-1$
\end{unitstatstable}

\unitimage{cerberus.png}

\unitdivider

\subsection{Falchion}

The Falchion uses a devastating weapon that combines the technologies employed by the Fellblade and Shadowsword super-heavy tanks. The Falchion's Twin Volcano Cannons are among the most powerful anti-tank weapons in the Imperium's arsenal. The Falchion requires such an investment of resources to build that its use is limited to the Legiones Astartes. Even then, it is so rare that most Legions keep only a handful, which they reserve for the destruction of the largest enemy war machines. As the Great Crusade spread ever further, the expeditionary fleets met an impressive number of enemies, but none could withstand the Falchion's formidable Volcano Cannon fire.

The Falchion has the \textbf{Wounds (3)} and \textbf{DR (5+)} abilities. Its Twin Volcano Cannons have the \textbf{Damage (+2)} ability.

\begin{unitstatstable}
\SetCell[r=3]{c,m} 15 cm &
\SetCell[r=3]{c,m} 2+ &
\SetCell[r=3]{c,m} +6 &
\SetCell[r=3]{c,m} 4 &
Twin Volcano Cannons &
90 cm &
1 &
2+ &
$-4$
\\
& & & &
Twin Heavy Bolters &
45 cm &
1 &
4+ &
$-1$
\\
& & & &
Laser Destroyer &
75 cm &
2 &
4+ &
$-3$
\end{unitstatstable}

\unitimage{falchion.png}

\unitdivider

\subsection{Fellblade}

The Fellblade is notable for its use of Mechanicum atomantic arc-reactor technology, combined with a reinforced chassis superior to that of the Baneblade. It is equipped with an Accelerator Cannon as well as secondary weapons such as a Demolisher Cannon and lascannons. This tank is extremely advanced, using many cutting-edge technologies. That is why its numbers are limited and its maintenance difficult. The Fellblade is crewed by Astartes. These war machines are used to break the most impenetrable enemy lines. They can also serve as siege tanks, siege-breakers, tank hunters, or mobile command fortresses.

The Fellblade has the \textbf{Wounds (3)} and \textbf{DR (5+)} abilities. Its Accelerator Cannon has the \textbf{Turret} and \textbf{Damage (+1)} abilities. Its Demolisher Cannon has the \textbf{Reduces Cover ($-3$)} ability.

\begin{unitstatstable}
\SetCell[r=3]{c,m} 15 cm &
\SetCell[r=3]{c,m} 2+ &
\SetCell[r=3]{c,m} +6 &
\SetCell[r=3]{c,m} 4 &
Accelerator Cannon &
75 cm &
2 &
3+ &
$-3$
\\
& & & &
Demolisher Cannon &
45 cm &
1 &
4+ &
$-3$
\\
& & & &
Laser Destroyer &
75 cm &
2 &
4+ &
$-3$
\end{unitstatstable}

\unitimage{fellblade.png}

\unitdivider

\subsection{Glaive}

The Glaive is a Fellblade variant deployed by the Space Marine Legions during the Great Crusade. It was designed by the Mechanicum around its main weapon, the Volkite Carronade, which is one of the most powerful weapons in a Legion's arsenal. This massive weapon incinerates a wide area in front of it, leaving only smoking slag. The Glaive was used actively by the 18 Legions, but some, such as the Dark Angels and the Salamanders, made more intensive use of it.

The Volkite Carronade hits every target in its line of sight along a line 3~cm wide and 45~cm long starting from its base.

The Glaive has the \textbf{Wounds (3)} and \textbf{DR (5+)} abilities. Its Volkite Carronade has the \textbf{Turret} and \textbf{Volkite Deflagration} abilities.

\begin{unitstatstable}
\SetCell[r=3]{c,m} 15 cm &
\SetCell[r=3]{c,m} 2+ &
\SetCell[r=3]{c,m} +6 &
\SetCell[r=3]{c,m} 4 &
Volkite Carronade &
30 cm &
Template &
3+ &
$-3$
\\
& & & &
Twin Heavy Bolters &
45 cm &
1 &
4+ &
$-1$
\\
& & & &
Laser Destroyer &
75 cm &
2 &
4+ &
$-3$
\end{unitstatstable}

\unitimage{glaive.png}

\unitdivider

\subsection{Kharybdis}

The Kharybdis assault claw can be considered an enlarged version of the Dreadclaw. It is a massive assault craft. The Kharybdis can transport a substantial strike force through the void while possessing enough firepower to cut a path through smaller vessels. The arrival of a swarm of these assault craft, boarding claws spread outward like hungry jaws, announces the death of the target upon which they fall. They are also used in ground assaults and, when they land, they unleash a swarm of fire that clears the surroundings. They can take off and move after landing. They also make combat use of their enormous thruster, burning enemies as they land or move. Their size is such that they can transport an entire formation or several Dreadnoughts.

The Kharybdis has the \textbf{Wounds (2)}, \textbf{Antigrav}, \textbf{Deep Strike (2)}, \textbf{Orbital Strike}, and \textbf{Transport (6)} abilities. Its Deathstorm uses a \textbf{Template (12 cm)} and has the \textbf{Bombing (1)} ability.

\begin{unitstatstable}
20 cm & 3+ & +0 & 4 & Deathstorm & Bomb & Template & 4+ & 0
\end{unitstatstable}

\unitimage{kharybdis.png}

\unitdivider

\subsection{Kratos}

The Kratos is based on an ancient Terran pattern used in large numbers during the Unification Wars. After the Treaty of Olympus, the Kratos was redesigned to serve in the Great Crusade. It serves as an infantry support tank, and its heavy guns quickly open a path for the infantry. The Kratos has nonetheless become obsolete and was withdrawn from service when the Legions began to favour faster and more versatile vehicles such as the ubiquitous Predator. Some have nevertheless been kept in the Legions' reserve armouries. In terms of firepower, the Kratos sits between the Sicaran and the Fellblade.

Both variants have the \textbf{Wounds (2)}, \textbf{DR (6+)}, and \textbf{Bulldozer} abilities. The Assault variant's Twin Multi-meltas have the \textbf{Turret} ability, and its Flamers have the \textbf{Reduces Cover ($-3$)} ability. The Support variant's Kratos Autocannon has two firing modes, and one must be chosen each turn. The Kratos Autocannon has the \textbf{Turret} ability, and the Kratos Autocannon -- High Explosive has the \textbf{Reduces Cover ($-1$)} ability.

\subsubsection*{Kratos -- Assault}

\begin{unitstatstable}
\SetCell[r=3]{c,m} 20 cm &
\SetCell[r=3]{c,m} 2+ &
\SetCell[r=3]{c,m} +3 &
\SetCell[r=3]{c,m} 4 &
Twin Multi-meltas &
20 cm &
2 &
3+ &
$-3$
\\
& & & &
Flamers &
20 cm &
3 &
4+ &
0
\\
& & & &
Autocannon &
45 cm &
2 &
4+ &
0
\end{unitstatstable}

\subsubsection*{Kratos -- Support}

\begin{unitstatstable}
\SetCell[r=4]{c,m} 20 cm &
\SetCell[r=4]{c,m} 2+ &
\SetCell[r=4]{c,m} +3 &
\SetCell[r=4]{c,m} 4 &
Kratos Autocannon -- High Explosive &
45 cm &
2 &
3+ &
0
\\
& & & &
Kratos Autocannon -- Armour-Piercing &
45 cm &
2 &
4+ &
$-2$
\\
& & & &
Heavy Bolter &
45 cm &
3 &
5+ &
$-1$
\\
& & & &
Autocannon &
45 cm &
2 &
4+ &
0
\end{unitstatstable}

\unitimage{kratos.png}

\unitdivider

\subsection{Spartan}

The Spartan was designed during the Great Crusade to transport a complete Terminator squad. Although it shares design features with the Land Raider, the Spartan is based on a more voluminous chassis. Its interior space is devoted entirely to its occupants, which makes it one of the terrestrial transports of choice. The Spartan has the advantage of being relatively fast for a vehicle of its size, which makes it even more effective in its role. Its armament is not neglected either, with numerous lascannons that allow it to threaten enemy armour.

The Spartan has the \textbf{Wounds (2)}, \textbf{All-Round Armour}, \textbf{Attached Transport}, and \textbf{Transport (4)} abilities.

\begin{unitstatstable}
20 cm & 2+ & +3 & Attached & Laser Destroyer & 75 cm & 2 & 4+ & $-3$
\end{unitstatstable}

\unitimage{spartan.png}

\unitdivider

\subsection{Thunderhawk}

A Thunderhawk can fulfil many tactical or strategic roles, though its primary mission remains troop transport; it excels in this role, able to carry several Space Marine squads if needed, as well as a substantial stock of ammunition and other equipment to support them in their task. They are often deployed from a planet's orbit from the armoured hangars of a cruiser. A squadron of Thunderhawks is enough to send an entire strike force into combat before the enemy even has the chance to realise the scale of the threat. Thanks to its multiple layers of armour, similar to those of Land Raiders, the Thunderhawk can withstand any atmospheric pressure and the fire of the heaviest anti-aircraft defences. Its powerful engines provide the energy needed to keep this armoured behemoth in the air as it dives roaring onto its targets. Astartes Assault Squads equipped with jump packs often deploy from the forward boarding ramp while the Thunderhawk is still moving, pouring a rain of fire onto their enemies and delivering death in the name of the Emperor.

Its Heavy Bolters have the \textbf{Turret} ability. The Thunderhawk has the \textbf{All-Round Armour}, \textbf{Wounds (2)}, \textbf{Flyer}, and \textbf{Transport (8)} abilities.

\begin{unitstatstable}
\SetCell[r=3]{c,m} 30 cm &
\SetCell[r=3]{c,m} 3+ &
\SetCell[r=3]{c,m} +1 &
\SetCell[r=3]{c,m} 4 &
Battle Cannon &
45 cm &
1 &
4+ &
$-2$
\\
& & & &
Lascannon &
45 cm &
2 &
5+ &
$-1$
\\
& & & &
Heavy Bolters &
30 cm &
3 &
5+ &
$-1$
\end{unitstatstable}

\unitimage{thunderhawk.png}

\unitdivider

\subsection{Thunderhawk Transporter}

A variant devoted to logistical support, the Thunderhawk Transporter allows a Legion's armour to be deployed rapidly onto the battlefield from orbit. It can also carry passengers. In a war zone, the Thunderhawk Transporter is armed with Heavy Bolter turrets for its own defence, but it is not intended to deliberately engage enemy forces. This pattern is used mainly during rapid strikes, or particularly intense campaigns, proving of inestimable strategic flexibility in the hands of a seasoned commander.

Its Heavy Bolters have the \textbf{Turret} ability. The Thunderhawk Transporter has the \textbf{All-Round Armour}, \textbf{Wounds (2)}, \textbf{Flyer}, and \textbf{Transport (4)} abilities, and may also transport 2 Vehicles.

\begin{unitstatstable}
30 cm & 3+ & +1 & 4 & Heavy Bolters & 30 cm & 3 & 5+ & $-1$
\end{unitstatstable}

\unitimage{thunderhawk-transporter.png}

\unitdivider

\subsection{Typhon}

The Typhon is named after the myths of ancient Terra and is armed with a massive cannon mounted on the chassis: the Dreadhammer siege cannon. This massive weapon was designed in part by Perturabo, Primarch of the Iron Warriors himself, and is a modified version of a static siege cannon that was used to reduce cities to dust. The weapon fires massive shells of several tonnes, powerful enough to liquefy flesh and bone and destroy even the most heavily fortified bunkers and strongpoints. Its structure is heavily reinforced to house the Dreadhammer Cannon, allowing it to pass without difficulty through the most obstructed areas. Apart from the tank's main weapon, the Typhon is only lightly armed compared with the Spartan or the Land Raider Proteus, with only two lascannons.

The Typhon has the \textbf{Wounds (2)}, \textbf{Bulldozer}, and \textbf{DR (6+)} abilities. Its Dreadhammer Cannon uses a \textbf{Template (7.5 cm)} and has the \textbf{Reduces Cover ($-3$)} ability.

\begin{unitstatstable}
\SetCell[r=2]{c,m} 15 cm &
\SetCell[r=2]{c,m} 2+ &
\SetCell[r=2]{c,m} +5 &
\SetCell[r=2]{c,m} 4 &
Dreadhammer Cannon &
75 cm &
Template &
3+ &
$-3$
\\
& & & &
Lascannon &
60 cm &
2 &
5+ &
$-1$
\end{unitstatstable}

\unitimage{typhon.png}
